\begin{statement}{C.2.D4.1}{Proved}{Global exclusion of connected positive-genus triple covers over a rational image}
Let $Y$ be a smooth compact complex threefold with
$b_2(Y;\Q)=0$, and let $f:Y\to B$ be a holomorphic map to a
smooth curve whose smooth nearby fibres are complex two-tori.
Fix a fibre $mS$, with $S$ reduced and irreducible, and restrict
the base to a small coordinate disc about it. Let
$H\subset S$ be a smooth rational curve along which the generic
transversal of $S$ is an ordinary triple point. Suppose that the
three generic normalization branches form one connected degree-three
cover $C\to H$, where $C$ denotes its smooth complete model.
Then $C$ is rational, and the cubic normal jet of $S$ vanishes
identically on at least one normal fibre over $H$.
In particular, a positive-genus connected degree-three cover is
impossible, and a primitive cubic normal jet is impossible for
any connected degree-three cover.
\end{statement}
\begin{proof}
We construct all comparison maps, including the map into the
actual nearby fibre. Put $N=N_{H/Y}$ and
$M=\OO_Y(S)|_H$. The ambient Chern-class hypothesis gives
$M\simeq\OO_H$ and $\deg N=-2$. The order-three normal jet is
therefore an actual homogeneous function
\[
 P\in H^0(H,\operatorname{Sym}^3N^*),\qquad
 U=\{v\in N:P(v)=1\}.
\tag{C.D4.1}
\]
The generic cubic has three distinct roots. The fibrewise Euler
identity $dP(v)(v)=3P(v)$ proves that $U$ is smooth everywhere,
including above a base point where the cubic degenerates. Above
a point where the entire cubic is zero there is no point of $U$.

\paragraph{Comparison with the actual fibre.}
Let $t$ be the base coordinate, so $f$ is a holomorphic function
on a neighbourhood of $H$. Locally $f=u s^m$, where $s=0$
defines $S$ and $u$ is a unit. The degree-$3m$ normal jet of
$f$ is $P^m$, after rescaling the chosen trivialization of $M$
by a constant. Indeed the ratio of the two jets is an invertible
holomorphic function on $H\simeq\PP^1$, hence constant.
Take a smooth tubular map $\tau$ from a neighbourhood of the
zero section in $N$ to $Y$, whose derivative on normal fibres
is the specified complex-linear identification. For positive
real $s\to0$, Taylor's formula gives
\[
 s^{-3m}f\bigl(\tau(sv)\bigr)\longrightarrow P(v)^m
                 \quad\hbox{in }C^\infty\hbox{ on compact subsets of }N.
\tag{C.D4.2}
\]
The order of vanishing is at least $3m$ along every point of
$H$: its ideal is a regular codimension-two ideal, and its
ordinary powers are primary along the smooth curve, so the
order verified at its generic point extends over $H$.

Choose any compact codimension-zero smooth submanifold
$K\subset U$. On a neighbourhood of $K$, the differential of
$P^m$ has a smooth right inverse as a real map to $\C$.
The implicit-function theorem, uniformly on $K$, or the
corresponding right-inverse vector field, transports this
neighbourhood to the level
$s^{-3m}f(\tau(sv))=1$. Composing with $v\mapsto\tau(sv)$
gives an orientation-preserving embedding into the actual
fibre $f^{-1}(s^{3m})$. In particular, intersections of compact
two-cycles, computed with parallel perturbations inside $K$,
are preserved. This construction uses the actual function
$f$ and its leading normal jet; it makes no assertion that
arbitrary changes of local equations preserve a Milnor fibre.

The algebraic surface $U$ admits a compact smooth core whose
inclusion is a homotopy equivalence. For example, resolve the
boundary in a projective compactification, triangulate the
resulting pair, take the complement of a sufficiently small
regular neighbourhood of the boundary and round its corners.
The radial collars of this regular neighbourhood give the
retraction. We choose this core for $K$. Thus the preceding
embedding induces one map carrying all of $H_1(U)$ and
$H_2(U)$, together with their compact intersection pairing.

Here is the precise torus obstruction we will use. For a map
$K\to T^4$, let $W$ be the rational image of $H_1(K)$ in
$H_1(T^4)$. The map lifts to the covering corresponding to
its fundamental-group image, a space homotopy equivalent to
$T^r$, where $r=\dim W$. Its image on $H_2$ consequently lies
in $\Lambda^2W\subset\Lambda^2\Q^4$. For $r\le3$ the wedge
pairing on this image is zero, because $\Lambda^4W=0$.
Consequently a compact core with $b_1(U)\le3$ and nonzero
intersection pairing cannot embed by (C.D4.2) into a nearby
torus. This is a comparison of the actual homology maps.

\paragraph{The normal bundle and the root divisor.}
Write
\[
                  N=\OO_H(a)\oplus\OO_H(-2-a),\qquad a\ge-1.
\tag{C.D4.3}
\]
If $a>0$, the coefficient of the cube of the first fibre
coordinate vanishes, since it has degree $-3a$. The other
coordinate is then a global linear factor of the cubic.
This contradicts connectedness of the degree-three root
cover. Thus $a=-1$ or $a=0$.

For $a=0$, irreducibility of the generic cubic forces the
coefficient of the first coordinate cubed to be a nonzero
constant. Along the trivial summand of $N$, the equation
$P=1$ therefore has three constant sections. Each is a
compact $\PP^1$ in $U$ with normal bundle $\OO_{\PP^1}(-2)$:
the derivative in the trivial summand is nonzero, so the
implicit-function theorem identifies its normal direction
with the other summand. Its homology square is $-2$.
Transporting its tubular neighbourhood by (C.D4.2) gives a
sphere of square $-2$ in a torus, whereas the torus has
$\pi_2=0$. This is impossible.

We are left with $N=\OO(-1)^{\oplus2}$ and
$E=\PP(N)\simeq\PP^1\times\PP^1$. Let $A$ be the common
zero divisor of the cubic coefficients on the base, of degree
$v$. Removing that scalar factor gives an integral horizontal
root divisor $R$ of bidegree $(3,3-v)$, where the first entry
is its degree over $H$. Its normalization is the given
cover $C$: at a general point the maps identify a normalization
branch with its distinct tangent line, and both complete
normal covers are the integral closure in this same function
field. Thus, writing $g=g(C)$,
\[
                       0\le g\le p_a(R)=4-2v.
\tag{C.D4.4}
\]
For a primitive cubic $v=0$. For a positive-genus cover with
$v>0$, the bound forces $v=1$. We treat these two cases.
Their total root divisors are respectively
\[
 B=R_{(3,3)},\qquad
 B=R_{(3,2)}+F_{(0,1)}.
\tag{C.D4.5}
\]
Both are reduced divisors of class $(3,3)$; the second has
one simple vertical component. They are connected. Let
$c=1$ or $2$ be their number of components, and set
$\nu=\sum_p(r_p(B)-1)$, where $r_p(B)$ is the number of
analytic branches at a singular point. Their total delta
invariants are $4-g$ and $5-g$, respectively, so in the
cases under consideration they are at most four.

\paragraph{The cyclic compactification and its local invariants.}
Put $L=\OO_E(1,1)$ and form the finite cyclic triple cover
\[
 X_0=\operatorname{Spec}_E
       (\OO_E\oplus L^{-1}\oplus L^{-2}),\qquad w^3=b,
\tag{C.D4.6}
\]
where $b$ is the section of $L^3$ defining $B$. The algebra
multiplication is given by $b$. Because $B$ is reduced,
$X_0$ is normal: its hypersurface local rings are Cohen--Macaulay
and regular in codimension one. The cover is connected,
as a meridian around any smooth point of $B$ acts by a
three-cycle. Above $E\setminus B$ it is exactly $U$.
Indeed a line $[v]\in E\setminus B$ has three radial vectors
with $P(v)=1$; equivalently, the dual tautological coordinate
has cube $b$. This gives the claimed isomorphism, including
the transition functions.

Let $\rho:X\to X_0$ resolve the surface singularities and
then the reduced boundary $D=X\setminus U$. Let $Q_p$ denote
the reduced exceptional divisor above a singular point of
$X_0$. We need the following finite computation:
\[
  \sum_p p_g(X_0,p)\le1,\qquad
  \sum_p b_1(Q_p;\Q)\le2.
\tag{C.D4.7}
\]
Here are its equations and derivation. A reduced plane germ
with delta invariant at most four has one of the forms
\[
\begin{array}{c|c|c}
 \hbox{type}&r(x,y)&\hbox{range}\\ \hline
 A_n&x^2+y^{n+1}&1\le n\le8\\
 D_n&x^2y+y^{n-1}&4\le n\le7\\
 E_6&x^3+y^4&\\
 E_7&x^3+xy^3&\\
 E_8&x^3+y^5&
\end{array}
\]
For completeness, multiplicity at least four gives
$\delta\ge\binom42=6$. At multiplicity two, Weierstrass
preparation and completing the square give $A_n$, with
$\delta=\lfloor(n+1)/2\rfloor$. At multiplicity three,
three different tangent lines give $D_4$. A double tangent
and a transverse line give that smooth line times an
$A_k$ branch or pair, adding intersection multiplicity two;
the delta bound gives $D_5,D_6,D_7$. With a single tangent
line, an irreducible branch has initial Puiseux pair $(3,4)$
or $(3,5)$, giving $E_6,E_8$; a smooth branch and a cusp
have total delta at least $1+3=4$, and equality gives $E_7$.
Three smooth branches with one tangent line already give
delta at least six. These reductions give the displayed
analytic forms; the higher terms in the last three forms
are eliminated by their positive weighted Jacobian terms.

The singularity of $X_0$ is $w^3=r(x,y)$. Normalize the
weights of $r$ to degree one and give $w$ weight $1/3$.
The adjunction calculation of
\cite[Theorem 2.1.1 and Corollary 2.1.3]{MerleTeissier1977}
counts $p_g$ by the positive lattice points under the
weight-one plane. Equivalently, it counts monomials
$x^iy^jw^k$ with
$(i+1)w_x+(j+1)w_y+(k+1)/3\le1$.
The normalized monodromy weights are those of the Jacobian
monomials, shifted by the sum of the variable weights
\cite[pp. 165--167]{Steenbrink1989Spectrum}.
The following table records the result; all entries are
also computed with rational arithmetic in
\path{code/root_cubic_normal_cone.py}.
\[
\begin{array}{c|c|c|c}
 \hbox{plane type}&\delta&p_g(w^3-r)&b_1(Q_p;\Q)\\ \hline
 A_1,A_2&1&0&0\\
 A_3,A_4&2&0&0\\
 A_5&3&1&2\\
 A_6&3&1&0\\
 A_7,A_8&4&1&0\\
 D_4&3&1&2\\
 D_5&3&1&0\\
 D_6,D_7&4&1&0\\
 E_6&3&1&0\\
 E_7,E_8&4&1&0
\end{array}
\]
To verify the last column directly, use the Jacobian bases
$\{y^j:0\le j<n\}$ for $A_n$;
$\{1,y,\ldots,y^{n-2},x\}$ for $D_n$;
$\{x^iy^j:0\le i<2,0\le j<3\}$ for $E_6$;
$\{1,y,\ldots,y^4,x,xy\}$ for $E_7$; and
$\{x^iy^j:0\le i<2,0\le j<4\}$ for $E_8$.
Count weights $(i+1)w_x+(j+1)w_y$ congruent to $1/3$ or
$2/3$ modulo one. The tensor-product monodromy for adding
$w^3$ identifies this count with the eigenvalue-one
multiplicity of the surface monodromy
\cite[\S1 and Example 2.4(a)]{Nemethi1991Join}.
The variation sequence of its simply connected Milnor fibre
identifies that multiplicity with the rational first Betti
number of its link. The negative-definite exceptional
intersection matrix then identifies $H_1$ of the link with
$H_1(Q_p)$: in the homology sequence of its plumbed
neighbourhood, the map on second homology is that invertible
matrix. This proves the column without replacing the actual
exceptional divisor by an unspecified graph.
Every nonzero entry in either last column costs delta at
least three. Since the total delta is at most four,
(C.D4.7) follows.

\paragraph{Homology and the compact intersection pairing.}
The finite algebra in (C.D4.6) and the K\"unneth formula give
$H^1(X_0,\OO_{X_0})=0$: the three summands on $E$ are
$\OO$, $\OO(-1,-1)$ and $\OO(-2,-2)$, all with $H^1=0$.
The resolution Leray sequence injects $H^1(X,\OO_X)$ into
$H^0(X_0,R^1\rho_*\OO_X)$, of dimension
$\sum p_g(X_0,p)$. Since $X$ is projective, (C.D4.7) gives
\[
                         b_1(X;\Q)=2q(X)\le2.
\tag{C.D4.8}
\]
The irreducible components of $D$ have independent homology
classes. The exceptional classes have a negative-definite
intersection matrix. In the first case of (C.D4.5) the
remaining class pushes forward to the nonzero class of $R$;
in the second, the two remaining classes push forward to
$(3,2)$ and $(0,1)$, which are independent on $E$.
Moreover their span contains the pullback of the ample
class $B=(3,3)$, whose square is positive. The Hodge index
theorem implies that the intersection form on this span is
nondegenerate: its part orthogonal to the positive pullback
class is negative definite.

Write $n_D$ for the number of boundary components. The exact
homology sequence of the pair $(X,U)$ identifies the map
$H_2(X;\Q)\to H_2(X,U;\Q)=\Q^{n_D}$ with intersection
against these components. Its surjectivity gives
$H_1(U;\Q)\simeq H_1(X;\Q)$. Its kernel is the orthogonal
complement of their span, and is the image of $H_2(U)$.
The pairing on that image is nondegenerate, of dimension
$b_2(X)-n_D$. As $D$ is connected,
$e(D)=n_D+1-b_1(D)$, and consequently
\[
\begin{split}
 b_1(U)&=b_1(X)\le2,\\
 \operatorname{rank}\bigl(\hbox{intersection pairing on }H_2(U)\bigr)
     &=b_2(X)-n_D\\
     &=e(U)-1+2b_1(X)-b_1(D).
\end{split}
\tag{C.D4.9}
\]
Here the equality of ranks follows also for a degenerate
pairing on $H_2(U)$: its kernel includes the kernel of the
map to $H_2(X)$, and the induced form on its image is the
nondegenerate form just identified.

There is a useful exact version allowing relations among
the strict boundary classes. If there are $\varepsilon$
independent relations, while their span still contains a
positive-square class, the same pair sequence and the same
Hodge-index argument give
\[
\begin{split}
 b_1(U)&=b_1(X)+\varepsilon,\\
 \operatorname{rank}\bigl(\hbox{intersection pairing on }H_2(U)\bigr)
 &=e(U)-1+2b_1(X)-b_1(D)+\varepsilon.
\end{split}
\tag{C.D4.10}
\]
These identities concern the actual classes in the indicated
cyclic cover. They will also be used for the degree-two
root divisor below.

To compute $b_1(D)$, start with the normalizations of the
$c$ components of $B$, of total genus $g$, and attach the
connected exceptional divisor over each singular point.
There are $r_p(B)$ attaching branches there. The incidence
graph contributes $\nu-c+1$ cycles. Subsequent point
blowups resolving the boundary subdivide edges or add trees
and preserve this first Betti number. Thus (C.D4.7) gives
\[
                         b_1(D)\le2g+\nu-c+3.
\tag{C.D4.11}
\]
The unramified cover $U\to E\setminus B$ has degree three,
so additivity and multiplicativity of Euler characteristic
give the exact value
\[
                         e(U)=3(4-2c+2g+\nu).
\tag{C.D4.12}
\]
For $c=1$, (C.D4.9)--(C.D4.12) give intersection rank at
least $3+4g+2\nu>0$. For $c=2$ they give rank at least
$4g+2\nu-2$. In this second case $g\ge1$, and $R$ meets
$F$ with intersection number three, so $\nu\ge1$ and the
rank is again positive. Both cases have $b_1(U)\le2$.
They contradict the compact-core comparison with $T^4$
proved immediately after (C.D4.2).

This excludes $v=0$ for every connected triple cover and
$v=1$ for every positive-genus one. By (C.D4.4) these are
all possibilities for a positive-genus cover, proving the
statement. The genus-zero cases with $v=1$ or $2$ are not
asserted to be excluded by this argument.
\end{proof}


% C.2.D4.3 from part_C2_cubic_cone.tex

